\section{Continuities with the Writings before Christ}\label{app:continuities}

Each teaching on the angels or the spirits that a pre-Christian writing
holds, set beside the Fathers who received, corrected, or rejected it and
the place where Aquinas or the Church takes the matter up. A parallel
without a naming or quotation by a Father is marked as a parallel.
Section~\ref{sec:before-christ} expounds each row.

\subsection*{Israel's writings: the Septuagint, the books outside the canon, and Philo}
\begingroup\footnotesize
\begin{longtable}{>{\raggedright\arraybackslash}p{0.17\linewidth}>{\raggedright\arraybackslash}p{0.22\linewidth}>{\raggedright\arraybackslash}p{0.25\linewidth}>{\raggedright\arraybackslash}p{0.14\linewidth}>{\raggedright\arraybackslash}p{0.14\linewidth}}
\toprule
\textbf{Witness and locus} & \textbf{Teaching} & \textbf{Patristic reception} & \textbf{Disposition} & \textbf{Aquinas or the Church}\\
\midrule
\endhead
\bottomrule
\endfoot
Septuagint, Deut 32:8--9 & Nations bounded by the number of the angels of God & 1 Clem. 29; Justin, Dial. 131; Clement, Strom. VII.2; Origen, De princ. I.5.2, C. Cels. V.29--30; Basil, Adv. Eun. III.1; Eusebius, Dem. ev. IV.7; Dionysius, CH 9.2--4 & Received & CCC 57 (cites Dt LXX 32:8); ST I q.108 a.6; q.113 a.8\\
Septuagint, Deut 32:8--9 & (Christological reading) nations pass from the angels to the Lord & Irenaeus, Adv. haer. III.12.9; Hilary, Tract. in Ps. 2.31 & Received (development) & ---\\
Septuagint, Job 1:6; 2:1 & Angels of God before the Lord, the devil among them & Gregory, Moralia II, as cited & Received & ST I q.112 a.3 ad 3\\
Septuagint, Job 38:7 & All the angels praised God when the stars were made & Augustine, De civ. Dei XI.9 (first day); Cassian, Conl. VIII.7 (before the world); Chrysostom, Hom. in Matt. 76 & Received; created with the corporeal world (common teaching) & ST I q.61 a.3; Lateran IV (DS 800)\\
Septuagint, Deut 32:43; Ps 96[97]:7 & Let all the angels of God worship him & Heb 1:6 & Received as Scripture & ---\\
Septuagint/Vulgate, Ps 137[138]:1 & Psalmody in the sight of the angels & Benedict, Regula 19 & Received & ---\\
Septuagint, Deut 33:2 & Angels at God's right hand at Sinai & Acts 7:53; Gal 3:19; Heb 2:2 & Received & ST I-II q.98 a.3\\
Septuagint, Isa 9:6 & Angel (Messenger) of great counsel & Justin, Dial. 126; Origen, C. Cels. V.53 & Received of Christ & ---\\
Septuagint, Isa 63:9 & Not an ambassador nor angel, but the Lord himself saves & Irenaeus, Adv. haer. III.20.4 & Received of the Incarnation & ---\\
Greek Daniel / Aramaic Daniel, Dan 3:49, 58; 4:10, 14 & Angel in the furnace; angels bless the Lord; the watchers & Augustine, De civ. Dei XI.9; Jerome, In Dan. 4:10 & Received as Scripture & Trent (DS 1502)\\
1 Enoch, 1:9 & The Lord comes with his holy ones & Jude 14--15; Tertullian, De cultu fem. I.3; Clement, Adumbr. in Iud.; Augustine, De civ. Dei XV.23 & Received through Jude; book Apocryphal & Trent (DS 1503); Jerome, De vir. ill. 4\\
1 Enoch, 6--8; 15--16 & Watchers took wives, begot giants; giants' spirits are demons; angels taught arts & Justin, 2 Apol. 5; Athenagoras, Leg. 24--25; Irenaeus IV.16.2, IV.36.4; Clement, Strom. V.1; Tertullian, De cultu fem. I.2, De idol. 9, De virg. vel. 7; Cyprian, De hab. virg. 14; Lactantius, Div. Inst. II.14; Sulpicius, Chron. I.2 & Received by these; rejected by Africanus (frag. 2), Chrysostom (Hom. in Gen. 22.2--3), Augustine (De civ. Dei XV.23; XVIII.38), Cassian (Conl. VIII.21) & ST I q.51 a.3 ad 6; q.63 aa.2, 4; Lateran IV (DS 800)\\
1 Enoch, 6:6 & Oath on Hermon & Hilary, Tract. in Ps. 132.6; Origen, Comm. in Io. VI.25 & Named and set aside & ---\\
1 Enoch, 10:4--12 & Fallen bound in darkness until the judgment & Jude 6; 2 Pet 2:4 (Augustine refers to first fall) & Received in substance & ST I q.64 a.4\\
1 Enoch, 19:1; 99:7 & Fallen spirits lead men to sacrifice to demons as gods & Tertullian, De idol. 4 (quotes as Enoch) & Received & Ps 95:5; 1 Cor 10:20\\
1 Enoch, 9:1--3; 40:6 & Archangels carry men's cause and prayers to God & Origen, C. Cels. V.4 (parallel) & Received & Tob 12:12; Apoc 8:3--4\\
1 Enoch, 14:22--23; 39:12; 71:7 & Those who sleep not sing the Sanctus & Jerome, In Dan. 4:10 (watchers, monastic vigil) & Received & Isa 6:3; Dan 7:10\\
1 Enoch, 61:10 & Angels of powers and of principalities with Cherubim, Seraphim & --- (parallel) & Received in substance & Col 1:16; ST I q.108 a.5\\
1 Enoch, 100:5 & Guardians from the holy angels for the righteous & --- (parallel) & Received & Matt 18:10; ST I q.113 a.2\\
1 Enoch; Tobit, 1 En 20; 40:9; Tob 12:15 & Seven archangels; offices of Raphael, Gabriel, Michael & Clement, Strom. VI.16; Origen, De princ. I.8.1 & Seven received; names beyond three rejected & Roman synod 745 (MGH Conc. II.1); Directory 217\\
Jubilees, 2:2--3 & Angels created on the first day, praising God's works & Epiphanius, De mensuris 22 (reproduces list, unnamed); Augustine, De civ. Dei XI.9 (parallel) & Parallel; created with the corporeal world (common teaching) & ST I q.61 a.3; Lateran IV ``simul'' (DS 800)\\
Jubilees, 2:2 & Angels over winds, clouds, seasons & Athenagoras, Leg. 24 (parallel) & Received & ST I q.110 a.1\\
Jubilees, 1:27--2:1 & Angel of the presence gives the Law's history & Acts 7:53; Gal 3:19 (parallel) & Received & ST I-II q.98 a.3\\
Jubilees, 4:15; 5:6--10 & Watchers first sent to instruct men, then bound & Justin, 2 Apol. 5; Lactantius, Div. Inst. II.14; Irenaeus IV.16.2 (parallels) & Corrected (sending of good angels kept; unions rejected) & ST I q.112\\
Jubilees, 10:8--11; 17:16; 48:10 & Tenth of spirits left to Mastema; evils permitted and limited & --- (parallel) & Received in substance & ST I q.114 a.1; q.64 a.4\\
Jubilees, 15:31--32 & Spirits over nations lead them astray; none over Israel & Dionysius, CH 9.3--4; Eusebius, Dem. ev. IV.7 & Corrected & CCC 57; ST I q.113 a.8\\
Jubilees, (book) & Little Genesis & Epiphanius, Pan. 39.6.1; Jerome, Ep. 78 & Apocryphal & ---\\
Testaments, T. Levi 3 & Ranked heavens; angels of the presence offer bloodless offering; thrones, dominions & Origen, C. Cels. V.4 (parallel) & Received in substance & Apoc 8:3; Col 1:16\\
Testaments, T. Levi 5; T. Dan 6 & Angel who intercedes for Israel, ``Mediator'' & Hermas, Sim. VIII.3 (Michael) & Corrected & 1 Tim 2:5; ST III q.26 a.1 ad 2\\
Testaments, T. Judah 20; T. Asher 6; T. Benj. 6 & Two spirits; angels of the Lord and of Satan; angel of peace guides the soul & Hermas, Mand. VI.2; Cassian, Conl. VIII.17 & Received & ST I q.113 aa.1--2; q.114 a.1\\
Testaments, T. Reuben 2--3 & Seven spirits of error & Origen, Hom. in Ies. Nave 15.5--6 (names the book) & Corrected & ST I q.114 a.3\\
Testaments, T. Reuben 5; T. Naph. 3 & Watchers allured by women & --- & Rejected with the Enochic reading & ST I q.51 a.3 ad 6\\
Philo, De gig. 16; De somn. I.141--142; De plant. 14 & Angels named by their office as messengers & Origen, C. Cels. V.4; Augustine, Enarr. in Ps. 103 s.1.15 & Received & ---\\
Philo, De gig. 6--16 & Angels, daemons, souls one kind & Origen, C. Cels. V.55 (recalls Philo's reading unnamed) & Rejected & Constantinople 543 (DS 403); ST I q.75 a.7\\
Philo, De conf. 171--175 & God rules through his powers as a king through ministers & --- (parallel) & Received and deepened & ST I q.103 a.6 and ad 3\\
Philo, De conf. 179; De opif. 72--75 & ``Let us make man'' said to subordinate powers & Justin, Dial. 62; Irenaeus, Adv. haer. IV.20.1 & Rejected & ST I q.45 a.5; q.91 a.4 ad 2\\
Philo, Quaest. in Gen. I.92 & Giants from angels and women & Ambrose, De Noe 4.8 & Received by Ambrose; rejected after Augustine & ST I q.51 a.3 ad 6\\
\end{longtable}
\endgroup

\subsection*{The Greek poets and philosophers}
\begingroup\footnotesize
\begin{longtable}{>{\raggedright\arraybackslash}p{0.17\linewidth}>{\raggedright\arraybackslash}p{0.22\linewidth}>{\raggedright\arraybackslash}p{0.25\linewidth}>{\raggedright\arraybackslash}p{0.14\linewidth}>{\raggedright\arraybackslash}p{0.14\linewidth}}
\toprule
\textbf{Witness and locus} & \textbf{Teaching} & \textbf{Patristic reception} & \textbf{Disposition} & \textbf{Aquinas or the Church}\\
\midrule
\endhead
\bottomrule
\endfoot
Hesiod, \emph{Works and Days} 121--126 & The golden race become pure spirits, guardians of mortal men & Lactantius, \emph{Div. inst.} II.14.7--8 (quotes 122--123) & Received: God sent guardians to men; corrected: the fallen usurp the guardians' name to be worshipped & \emph{ST} I q.113 a.2; Lateran IV, DS 800 (angels created from nothing)\\
Hesiod, \emph{Works and Days} 252--255 & Thrice ten thousand watchers of mortal men & Eusebius, \emph{Praep. ev.} IV.17 & Rejected as applied to the pagan gods and good daemons, who did not guard & Parallel: Dan 4:10, 14 (``watcher'')\\
Plato, \emph{Cratylus} 397e--398c & Daemon named from ``knowing''; the good man becomes a daemon & Lactantius, \emph{Div. inst.} II.14.6; Augustine, \emph{De civ. Dei} IX.20; Isidore, \emph{Etym.} VIII.11.15 & Etymology received; corrected: knowledge without charity puffs up; the name belongs to the fallen & Aquinas, \emph{De sub. sep.} 20\\
Plato, \emph{Apology} 31c--d, 40a--c & Socrates' sign from childhood, which forbids & Justin, \emph{1 Apol.} 5, \emph{2 Apol.} 10; Tertullian, \emph{Apol.} 22, \emph{De anima} 1; Minucius, \emph{Oct.} 26; Lactantius, II.14.9; Clement, \emph{Strom.} V.14 & Disputed among the Fathers: Socrates a foe of the demons (Justin); his spirit an evil demon (Tertullian, Minucius, Lactantius); named without censure (Clement) & ---\\
Plutarch, \emph{De genio Socratis} 20 & The daemon's thought imprinted on the mind as light, without voice & --- & Parallel & \emph{ST} I q.111 a.1 (angelic enlightenment of man)\\
Plato, \emph{Symposium} 202d--203a & The daemonic interprets between gods and men, carrying prayers up and commands down & Minucius, \emph{Oct.} 26; Lactantius, II.14.9; Origen, \emph{C. Cels.} V.4--5 & Office received and given to the angels, named from their office; worship reserved to God through Christ; ``demon'' reserved to the wicked & \emph{ST} I q.64 a.4 (angels midway between God and men); CCC 329, 331\\
Plato (via Apuleius), \emph{Symposium} 203a; Apuleius, \emph{De deo Socr.} 4 & ``God mingles not with man'' & Augustine, \emph{De civ. Dei} VIII.18, 20; IX.15 & Rejected: Christ the one Mediator; demons mediate only misery & CCC 331 (Christ the centre of the angelic world)\\
Plato, \emph{Republic} X 617d--e, 620d--e & A daemon allotted to each soul, guardian of its life & Clement, \emph{Strom.} V.14 & Received as a figure of the tutelary angels (Matt 18:10) & \emph{ST} I q.113 aa.2, 5\\
Plato, \emph{Republic} X 617e & ``The blame is his who chooses, God is blameless'' & Justin, \emph{1 Apol.} 44 & Received as borrowed from Moses & ---\\
Plato, \emph{Phaedo} 107d--108b & Each man's genius leads his soul to judgment & Tertullian, \emph{De anima} 53 & Corrected: the angel who receives the soul, ``the Mercury of the poets''; parallel Luke 16:22 & ---\\
Plato; Apuleius, \emph{Republic} X 620d--e; \emph{De deo Socr.} 15--16 & The genius or allotted daemon given at birth & Tertullian, \emph{De anima} 39; Origen, \emph{C. Cels.} VIII.34; Lactantius, II.14.12 & Pagan genius rejected as a demon (Tertullian, Lactantius); guardian angel of the little ones affirmed (Origen) & \emph{ST} I q.113 a.5 (guardian from birth)\\
Plato, \emph{Timaeus} 40d & The origin of the other daemons is beyond telling & Athenagoras, \emph{Leg.} 23; Eusebius, \emph{Praep. ev.} XIII.1 & Read as Plato's veiled rejection of the theogonies & ---\\
Plato, \emph{Timaeus} 41a--b & Made gods, immortal only by the Maker's will & Augustine, \emph{De civ. Dei} IX.23 & Received as the holy angels, if created and blessed by cleaving to God & \emph{ST} I q.50 a.5 obj. 2 and ad 2 (Aquinas reads them as the heavenly bodies; angels incorruptible by nature)\\
Platonists, (as reported by Eusebius) & Rational natures and daemons as effluences, not made from nothing & Eusebius, \emph{Praep. ev.} XIII.15 & Corrected: created from nothing; demons from the fall of angels & Lateran IV, DS 800\\
Plato, \emph{Laws} IV 713c--e; \emph{Statesman} 271d & Daemons set by God as shepherds over men and kinds of living things & Athenagoras, \emph{Leg.} 24 (parallel: angels' particular providence) & Received & \emph{ST} I q.110 a.1 ad 3 (``The holy doctors held with the Platonists'')\\
Plato (via Nemesius), threefold providence & Supreme God, heavenly gods, and demons divide providence & --- (Aquinas) & Rejected as to the design of government; intermediaries received as executors & \emph{ST} I q.22 a.3; q.103 a.6 ad 1\\
Plato, \emph{Laws} X 896e, 906a & A good and an evil soul; gods and daemons our allies in an immortal conflict & Clement, \emph{Strom.} V.14; Eusebius, \emph{Praep. ev.} XI.26 & Received as from Moses: the devil, the spiritual combat (Eph 6:12), Job 1:6, Deut 32:8 LXX & ---\\
Aristotle, \emph{Metaphysics} XII.8, 1073a--1074a & Unmoved movers as many as the heavenly motions, by probable argument & --- (Aquinas) & Corrected: the angels exceed all material multitude & \emph{ST} I q.50 a.3; \emph{SCG} II.92; \emph{De sub. sep.} 2\\
Aristotle, \emph{Metaphysics} XII.8, 1074b & The first substances are gods, an ``inspired utterance'' & Aquinas, \emph{In Metaph.} XII lect. 10 n. 31 & Received as to immaterial substances; one God only & ---\\
Aristotle, \emph{Metaphysics} XII.8 & Spiritual substances move only the heavenly bodies; no demons & --- (Aquinas) & Corrected: angels govern the lower bodies too; demons explain possession and magic & \emph{ST} I q.110 a.1 ad 2; \emph{De sub. sep.} 2, 4\\
Plutarch, \emph{De defectu oraculorum} 10--17 & Daemons midway, officers of the gods, some evil; oracles fail when they depart; ``the great God Pan is dead'' & Eusebius, \emph{Praep. ev.} V.1, 3--5, 16--17 & Rejected as to worship: the oracles were demons', silenced under Tiberius by Christ & \emph{ST} II-II q.94 a.4 and ad 2\\
Apuleius, \emph{De deo Socratis} 6, 13 & Aerial, passive, eternal daemons carry prayers and gifts & Augustine, \emph{De civ. Dei} VIII.14--22; IX.19 & Rejected: the air is the demons' prison; good spirits are angels, never ``good demons'' & \emph{ST} I q.51 a.1 ad 1; q.64 a.4\\
Platonists (via Augustine), gods above the moon, demons below & Division of the heavens & Augustine (as cited in \emph{ST}) & Received at one point & \emph{ST} I q.63 a.7 (with Damascene as cited there)\\
Plotinus, \emph{Enneads} III.4.3 & The guardian is the next higher power of one's own soul & --- & Parallel and contrast: the guardian angel is a person & \emph{ST} I q.113 aa.2, 5\\
Plotinus, (as reported by Augustine) & Heavenly spirits and men receive light from one Light, God & Augustine, \emph{De civ. Dei} X.2 & Received (John 1:9) & ---\\
Porphyry, \emph{De abstinentia} II.38 & Good daemons over animals, fruits, rains, winds; the ``transporters'' of the \emph{Symposium} & Origen, \emph{C. Cels.} VIII.31--32 (the same doctrine in Celsus) & Received as the angels, ``invisible husbandmen and guardians'' & \emph{ST} I q.110 a.1 ad 3\\
Porphyry, \emph{De abstinentia} II.38--42 & Evil daemons crave worship, sacrifice, and the name of the supreme god & Eusebius, \emph{Praep. ev.} IV.22 & Received as the pagans' own witness against their cult & \emph{ST} II-II q.94 a.4\\
Porphyry, theurgy; angels declaring the Father & Theurgic purification; angels who reveal to theurgists & Augustine, \emph{De civ. Dei} X.9--11, 26, 29 & Theurgy rejected; angels to be imitated, not invoked; the way is the Incarnation & ---\\
Iamblichus, \emph{De mysteriis} II.3 & Gods, archangels, angels, daemons, archons, souls known by theurgic apparitions & Augustine, \emph{De civ. Dei} X.9--10 & Rejected: demons invoked under the names of angels (2 Cor 11:14) & ---\\
Proclus (via \emph{Liber de causis}); Dionysian corpus, \emph{In De causis} prooem.; \emph{In De div. nom.} prooem. & Order of first causes; Platonic style of Dionysius & Aquinas, \emph{In De causis} prooem.; \emph{In DN} prooem.; \emph{De sub. sep.} 18 & Received as to the first principle; separate forms rejected; all spirits created immediately by God & \emph{ST} I q.112 a.4 ad 2\\
Greek philosophy (general), --- & Truths about God and the spirits & Justin, \emph{1 Apol.} 44, 59--60, \emph{2 Apol.} 10, 13; Clement, \emph{Strom.} I.5, I.17, I.21, V.14, VI.17, VII.2; Eusebius, \emph{Praep. ev.} XIII pref., 13--14; Augustine, \emph{De civ. Dei} VIII.5, 9, 11, \emph{De doctr. chr.} II.40.60 & Received as seeds of the Word, borrowings from Moses, gifts through the angels of the nations, the gold of Egypt; corrected where they contradict & \emph{ST} II-II q.94 a.1 (``we call them angels''); I q.110 a.1 ad 3\\
Hebrew doctrine set against Greek daemonism, --- & Rational powers ordered like stars; the fallen; ``the gods of the nations are daemons'' & Eusebius, \emph{Praep. ev.} VII.15--16 & Hebrew doctrine preferred to Greek polytheism & Ps 95[96]:5\\
\end{longtable}
\endgroup
